\chapter{三平均：最小奇素数的完整证明}\label{ch:three}

\section{三平均需要补哪些位置数}
奇数下半带已经覆盖 $p=3,n=7$；九位置可排成 $3\times3$，先平均各行再平均各列。
因此还需 $n=8,10$，以及全部 $n\ge11$。
这里采用已有的三平均直接证明，不强行把 $p\ge5$ 的容量公式应用到 $p=3$。

大维数仍用两个三重值与整数能量下降。
八元在一个 $(4,3,1)$ 核心上作欧几里得式下降。
十元在一个 $(6,3,1)$ 核心上比较七个返回对应的二次势能。
后者的有限信息是十二行二次多项式表，而不是对无限多个有理输入逐一搜索。

\section{全部n至少11：双三重值的整数下降}

\begin{lemma}[三平均的全部n至少11]\label{lem:triple-large}
三平均在全部 $n\ge11$ 上满足 $G$ 为三幂的充要条件。
\end{lemma}

\begin{proof}
\par\noindent\textbf{准备：统一的保护容量。}
记 $d=d_3(n)$，$f(a)$ 表示实际数值 $a$ 的重数。
有
\[
n\ge d+10.
\]
当 $d\le1$ 时显然；$d=2$ 时 $n\ge2\cdot7=14$；
当 $d\ge3$ 时，非三素因子的最小乘积给 $n\ge2\cdot5^{d-1}\ge d+10$。
每个危险主类的补集有两或三个位置，预留至多 $d$ 个位置便能同时保护。
在单一模三余数类的轻值池中，每个实际值至多两份：
池大小达到 $d+3$ 时能安全地选一个非恒等整数三元组，达到 $d+4$ 时还能避开一个指定均值。
后一结论使用四个候选数的删一比较。

从 $3X$ 开始。若已有两个三重值，无须初始化。
只有一个三重值 $a$ 且 $f(a)\ge5$ 时，平均 $(a,a,b)$，$b\ne a$，
留下三份 $a$ 并产生新的三重值。
若 $f(a)=3,4$，其外至少 $n-4\ge d+6$ 个位置，可用保护池选出均值不同于 $a$ 的三元组。
若没有三重值，先用保护池造出一个 $a$；它至多有五份，
其外至少 $d+5$ 个未改变的初值仍被 $3$ 整除，第二步再造一个不同三重值。

下面给每个非终端状态一个保持两个三重值的安全整数续步。
先说明一种将反复用到的安全性：若未选位置留下 $a,b$，而 $x=3b-2a$，
则假设输出全同余于 $c$ 会先得到 $a=b=c$，再得到 $x=c$。
因此选择 $(a,x,x)$ 时，即使危险素数为 $2$，也无须除以 $2$ 来证明安全。
同理，$x=3b-2a,y=3a-2b$ 时，$(x,y,z)$ 和 $(x,b,z)$ 均可用未选锚点及组和证明安全。

\par\noindent\textbf{第一类：至少三个三重值。}
若至少有四个，模三只有三个类，故有两个不同重值 $a,b$ 同余。
平均 $(a,a,b)$，至少另外两个三重值未动。
若恰三个，且 $a,b$ 同余，记第三个为 $c$。
两个候选 $(a,a,b)$、$(b,b,a)$ 的均值不同，所以至少一个不同于 $c$，保留 $c$ 与新均值即可。

若三个重值分处三个模三类，且另有轻值 $x$，取同余于 $x$ 的重值 $a$，平均 $(a,a,x)$，
另外两个重值保留。
若只有三个实际值，平均 $(a,b,c)$；若至少两个值有四份，它们仍有三份。
只剩重数为 $3,3,n-6$ 的情形，设最后一个值为 $c$。
如果新均值碰撞到 $c$，便有 $a+b=2c$；零和进一步给 $nc=0$。
于是得到明确终端
\[
(a^3,(-a)^3,0^{n-6}),
\]
三次 $(a,-a,0)$ 即可完成。

\par\noindent\textbf{第二类：两个三重值模三同余。}
设它们为 $a,b$。若 $f(a)\ge4$，用 $(b,b,a)$；若 $f(b)\ge4$，对称处理。
至少一个旧值仍有三份，新均值与它不同。
因此只剩 $f(a)=f(b)=3$，其外的轻值池 $S$ 满足 $|S|=n-6\ge d+4$。

若 $S$ 中有与锚点同余的 $x$，两个候选 $(a,a,x)$、$(b,b,x)$ 都安全。
若都碰撞，必有 $2a+x=3b$、$2b+x=3a$，相减得到 $5(a-b)=0$，矛盾。
若 $S$ 只占另外一个余数类，直接用保护池。
若它占另外两个类 $S_1,S_2$，考虑跨类对 $(c,e)$。
候选 $(a,c,e)$ 和 $(b,c,e)$ 不能选中的危险，仅可能是某个二元危险集恰等于 $\{c,e\}$。
每个危险素数至多阻塞一个跨类位置对，而
\[
|S_1||S_2|\ge |S|-1=n-7>d.
\]
故有未被阻塞的对。两个安全候选若同时碰撞，会给
$a+c+e=3b$、$b+c+e=3a$，推出 $4(a-b)=0$，仍矛盾。

\par\noindent\textbf{第三类：两个三重值模三不同。}
若轻值 $x\equiv a\pmod3$，先试 $(a,a,x)$。
失败只可能是 $f(a)\le4$ 且 $x=3b-2a$。
若这个 $x$ 出现两份，改用 $(a,x,x)$，得到 $2b-a$，与 $a,b$ 均不同。
其安全性由前述仿射见证保证。
因此还未成功时，$a$ 类至多有一个例外单点 $x=3b-2a$，
$b$ 类至多有一个例外单点 $y=3a-2b$，其余位置都在第三类，记其池为 $Z$。

若 $x,y$ 都存在，$f(a)+f(b)+2\le10<n$，所以 $Z$ 非空。
任取 $z\in Z$，平均 $(x,y,z)$，两个锚点完整保留，且 $x-y=5(b-a)\ne0$。

若只有 $x$，先说明 $Z$ 非空。否则安全操作 $(a,a,x)$ 后只有两个实际值 $a,b$，
而它们模三不同。合法性使 $a-b$ 与 $n$ 的每个非三素因子互素，
模三不同又补上素数 $3$，故 $\gcd(a-b,n)=1$。
零和迫使 $n\mid f(a)$，与 $0<f(a)<n$ 矛盾。
所以可选 $(x,b,z)$。若 $f(b)\ge4$，两个旧重值均保留。
若 $f(b)=3$，则 $|Z|\ge n-8\ge3$，而每个轻值至多两份，
可以选 $z\ne5a-4b$，使新均值不同于保留的 $a$。
只有 $y$ 的情形对称。

若 $x,y$ 都不存在，$Z$ 仍不能为空，理由同上。
当 $f(a)=f(b)=3$ 时，用 $|Z|\ge d+4$ 的保护池；
当二者均至少四份时，任取 $(a,b,z)$，两个旧重值仍保留。
最后设 $f(a)=3,f(b)\ge4$，选 $z\ne2b-a$ 即可。
若无此选择，$Z$ 全部等于 $2b-a$，其位置数 $r$ 只能为 $1$ 或 $2$。
零和变成
\[
(3-r)(a-b)+nb=0.
\]
所有值都在 $b+(a-b)\Z$ 中，所以合法性与模三不同仍给 $\gcd(a-b,n)=1$。
上式却迫使 $n\mid3-r$，不可能。

\par\noindent\textbf{有限终止。}
三类已经穷尽全部储备状态。每个非终端操作都安全、非恒等，并保持整数与两个不同三重值。
平方和严格下降且为非负整数，故有限步后必到前述相反值终端。
对 $3X$ 得到的同一平均词也使 $X$ 归零。
\end{proof}



\section{八位置：将第二坐标的绝对值降到一}
八元先用安全整数平均制造一个五重值，随后压成 $B_8(u,v)=(u^4,v^3,-4u-3v)$。
这一步以后改用参数下降，不再要求实际中间数全是整数。
取本原方向 $(x,y)=(u,u+v)$，合法性使 $y$ 为奇数。
如果能用真实返回把 $|y|$ 严格降低，那么它作为正整数也会有限终止。

\begin{lemma}[三平均八元的欧几里得下降]\label{lem:eight}
三平均八元的合法配置均可归零。
\end{lemma}
\begin{proof}
若无五重值，存在安全非恒等整数三元组。奇偶分布为 $4+4$ 或 $2+6$；前者任意三元组补集都保留奇偶见证。若没有非恒等整三元组，则至多两个模三类，且每个至少三项的类实际常值；无五重值迫使 $(a^4,b^4)$，由零和得 $b=-a$，违反合法性。在 $2+6$ 情形，若多数池没有整三元组，无五重值迫使它的两模三类各至少两项；其两项和覆盖全部模三类，任一少数位置可与其中两项配成安全整三元组。
整数能量下降至五重值 $u$ 后，写例外为 $a,b,c$，选 $c-u$ 奇。平均 $(u,a,b)$，得到合法 $B_8(u,v)=(u^4,v^3,-4u-3v)$。

\par\noindent\textbf{在核心上作下降。}
令 $x=u,y=u+v$，本原参数满足 $y$ 奇。下列实际返回给三个作用：
\[
R:(x,y)\mapsto(x+y,-y),\quad
S:(x,y)\mapsto(y-x,y),\quad
F:(x,y)\mapsto(3x+3y,2x-4y).
\]
为列出操作，置 $w=-4u-3v$、$a_0=(u+2v)/3$、$b_0=(2u+v)/3$、$q_0=-u-2v/3$、$h_0=-u/3-4v/9$、$k_0=u/3-2v/9$。三个词依次为
\begin{align*}
R:&\ (v,v,u),(v,u,u),(w,a_0,b_0),(a_0,b_0,u);\\
S:&\ (w,v,u),(q_0,q_0,u),(q_0,u,u),
(h_0,v,k_0),(h_0,v,k_0),(v/9,v/9,k_0);\\
F:&\ (v,u,u),(w,b_0,b_0),(v,u,u).
\end{align*}

这里每个括号表示在\emph{当时}选取的三个值。
新值 $a_0,b_0,\ldots$ 是前面步骤产生的值，反复写同一值时要取不同位置。
这些词的输出重新具有重数 $4,3,1$，并在整体非零尺度下给出所列参数矩阵。
因此它们可以在任意后继核心上重放。
逐次代入可核对库存及最终重数 $4,3,1$。$SR,RS$ 给平移 $x\mapsto x\pm2y$。若 $|y|>1$，先用 $S$ 使 $x$ 奇，再平移至 $|x_0-2y|<|y|$；等号仅在 $y\mid x$ 时发生，与本原性矛盾。
余数的平衡选择与通常的欧几里得算法相同：允许 $x$ 加减 $2y$ 后，
把 $x_0$ 选在以 $2y$ 为中心、宽度为 $2|y|$ 的区间内。
边界等号意味着 $y\mid x$，本原性会迫使 $|y|=1$，与当前阶段矛盾。
施加 $F$ 后两坐标均偶，约去 $2$ 后第二坐标奇，故本原输出满足 $|y'|\le|x_0-2y|<|y|$。最终 $|y|=1$，平移和 $S$ 送到 $(0,\pm1)$。配置 $(0^4,1^3,-3)$ 先平均 $(-3,0,0)$，再三次平均 $(-1,0,1)$。
\end{proof}


\section{十位置：先归约，再用二次势能下降}
十元的危险素数为 $2$ 和 $5$。
入口阶段用未选位置保留这两种见证，先制造六重值，
再用四步桥接进入 $(6,3,1)$ 核心。
进入核心后，整数势能由本原参数及其模 $4$ 类型决定。
下面先说明无限输入怎样归约到这个势能，随后给出全部七个返回和十二行有限表。

\begin{lemma}[三平均十元的二次势能下降]\label{lem:ten}
三平均十元的合法配置均可归零。
\end{lemma}
\begin{proof}

\par\noindent\textbf{第一步：取得六重值。}
先证明无六重值时有安全整三元组。危险奇偶集只能有两项，记为 $P$；危险模五集有两或三项，记为 $F$。若二者均不存在，任意整三元组安全；不存在整三元组会迫使 $(a^5,b^5)$，违反合法性。仅有一种危险集时，在多数池内用模三两类的两项和；若无法配出安全三元组，多数池中至少六项为同一实际值。
两危险集同时存在时，令 $K=(P\cup F)^c$。其内部整三元组安全。若 $|K|\ge7$，无六重值使两个模三类均至少两项，与任一外部位置配对即得安全选择。$|K|=6$ 时只剩 $r^5s$，全部外部余数必须为 $s$；从 $P\setminus F,F\setminus P$ 各取一项，连同 $K$ 中的 $s$ 即可。$|K|=5$ 时 $P,F$ 不交，只剩 $r^4s$ 或 $r^5$；前者同理。后者若 $P,F$ 有跨类余数，配一项 $K$ 得三类各一；否则外部五项同余，取一项 $P$、两项 $F$。每种选择补集均保留两种见证。


\par\noindent\textbf{第二步：从六重值到二维核心。}
能量下降至 $(u^6,a,b,c,d)$。令 $e_a=a-u$ 等，和为 $-10u$。若某个例外差是模 $10$ 单位，平均其余三个例外，一步进入合法 $B_{10}(u,v)=(u^6,v^3,-6u-3v)$。否则四个差模 $10$ 为两份 $5$ 及一对相反非零偶数，故有单位和的跨类配对，设为 $e_a+e_b$。令
\[q=(-5u-a-b)/3,\quad r=(-5u+2a+2b)/9,\quad s=(-11u-a-b)/27.
\]
词 $(u,c,d),(q,a,b),(u,q,r),(u,q,r)$ 给 $u^3r s^6=B_{10}(s,u)$，且 $s-u=-(40u+e_a+e_b)/27$ 是模 $10$ 单位。


\par\noindent\textbf{第三步：用有限局部类型选择下降返回。}
在本原 $(X,Y)$ 方向 $X=u,Y=3u+2v$ 上，合法类型为：两者奇且 $X+Y\equiv0\pmod4$，取 $h=1$；或奇偶相反，取 $h=2$。另要求 $5\nmid Y$。整数势能为 $h^2(5X^2+Y^2)$。命题\ref{prop:ten-certificate}给七个实际返回及十二个完整二次覆盖：除 $X=0$ 外，总有返回严格降低该势能。$X=0$ 时已有零坐标，固定一个零，对其余九个位置采用两层三进网络。由整数下降完成。
\end{proof}



\section{七个返回与全部十二个二次覆盖}
下面的有限命题正是十元证明所用的有限输入。
矩阵 $M_i$ 给参数的线性变化；本原化内容 $g_i$ 和输出类型 $h_i$ 只依赖有限局部信息。
对于固定类型，把 $t=Y/X$ 视为实变量，
“势能严格下降”便成为一个显式二次多项式为正。
某一行的多个多项式覆盖整个允许方向，就意味着这一行中的任意真实有理输入总有一个可选返回。

表中列出每一种类型的全部二次系数及剩余方向。
无穷方向 $t=\infty$ 对应 $X=0$，本来就是已知终端。
有限方向 $t=0$ 对应 $Y=0$，由合法性 $5\nmid Y$ 排除。

\begin{proposition}[十元二次覆盖证书]\label{prop:ten-certificate}
引理\ref{lem:ten}的合法本原方向上，以下七个矩阵中至少一个严格降低 $h^2(5X^2+Y^2)$，唯一相容例外为 $X=0$：
\[
\begin{gathered}
M_1=\begin{pmatrix}1&1\\-5&-1\end{pmatrix},\quad
M_2=\begin{pmatrix}1&0\\-5&2\end{pmatrix},\quad
M_3=\begin{pmatrix}1&-2\\0&3\end{pmatrix},\quad
M_4=\begin{pmatrix}1&1\\5&-3\end{pmatrix},\\
M_5=\begin{pmatrix}4&1\\0&-3\end{pmatrix},\quad
M_6=\begin{pmatrix}1&-2\\5&2\end{pmatrix},\quad
M_7=\begin{pmatrix}1&0\\-10&-1\end{pmatrix}.
\end{gathered}
\]
\end{proposition}
\begin{proof}
置 $U=u,V=v,W=-6u-3v$，$Q=(-5u-2v)/3,R=(u+2v)/3,S=(2u+v)/3,T=-2u-v/3$。返回依次平均的当前三元组为
\begin{center}\small
\begin{tabular}{cl}\toprule
矩阵&平均词\\\midrule
$M_1$&$(W,V,U),(V,U,U),(V,U,U)$\\
$M_2$&$(W,V,U),(Q,U,U),(V,V,U),(Q,R,U),(Q,R,U)$\\
$M_3$&$(W,U,U),(V,U,U),(V,S,U),(V,S,U)$\\
$M_4$&$(W,U,U),(V,U,U),(V,U,U)$\\
$M_5$&$(W,V,U),(V,V,U),(R,U,U),(R,U,U)$\\
$M_6$&$(V,U,U),(W,S,U),(V,S,U),(V,S,U)$\\
$M_7$&$(W,V,V),(V,U,U),(T,S,U),(T,S,U),(S,U,U)$\\\bottomrule
\end{tabular}\end{center}
代入两个独立参数可核对全部线性形式和库存。行列式为 $4,2,3,-8,-12,12,-1$，所以本原化内容只含 $2,3$。模 $96$ 记录各内容及输出模四类型，得到5120个合法局部类和12种签名。
若 $M_i=(a,b;c,d)$、内容为 $g_i$、输出尺度为 $h_i$，
输出本原列是 $M_i(X,Y)^T/g_i$，所以输出势能为
\[
\frac{h_i^2}{g_i^2}
\left(5(aX+bY)^2+(cX+dY)^2\right).
\]
把它与输入势能比较，乘正数 $g_i^2$，再在 $X\ne0$ 时除以 $X^2$，
下降条件是
\[
g_i^2h^2(5+t^2)-h_i^2[5(a+bt)^2+(c+dt)^2]>0,\qquad t=Y/X.
\]
表\ref{tab:ten-data}列出所有系数。每行各多项式的实根把实直线分成有限区间，逐区间及逐根检查符号；二次项检查无穷方向。剩余的有限点只可能为 $0$，被 $5\nmid Y$ 排除；两行留下无穷方向，其中仅一行与 $(0,1)$ 的局部签名相容，正是终端 $X=0$。由此得到完整覆盖。
\end{proof}

\begingroup\small
\begin{longtable}{r r l p{83mm} p{24mm}}
\caption{十二个二次覆盖：$i:(a,b,c)$ 表示 $a+bt+ct^2>0$}\label{tab:ten-data}\\
\toprule 类&个数&代表&$i:(a,b,c)$&未严格覆盖\\\midrule\endfirsthead
\toprule 类&个数&代表&$i:(a,b,c)$&未严格覆盖\\\midrule\endhead
1&768&$(0,1)$&1:(-10,-20,-2); 2:(-40,80,0); 3:(0,80,-112); 6:(50,0,-8); 7:(-400,-80,0)&$\infty$\\
2&768&$(1,0)$&1:(-10,-20,-2); 2:(-10,20,0); 3:(0,80,-112); 5:(0,-160,8); 7:(-400,-80,0)&0\\
3&256&$(1,2)$&1:(-10,-20,-2); 2:(-10,20,0); 3:(160,80,-80); 5:(640,-40,130); 7:(-400,-80,0)&$\varnothing$\\
4&384&$(1,3)$&1:(-40,-80,-8); 3:(0,20,-28); 4:(50,20,2); 7:(-100,-20,0)&$\varnothing$\\
5&768&$(1,6)$&1:(-10,-20,-2); 2:(-10,20,0); 3:(0,80,-112); 5:(0,-40,2); 7:(-400,-80,0)&0\\
6&384&$(1,7)$&1:(-40,-80,-8); 3:(0,20,-28); 4:(200,80,8); 7:(-100,-20,0)&$\varnothing$\\
7&256&$(1,8)$&1:(-10,-20,-2); 2:(-10,20,0); 3:(160,80,-80); 5:(2560,-160,520); 7:(-400,-80,0)&$\varnothing$\\
8&128&$(1,11)$&1:(-40,-80,-8); 3:(40,20,-20); 4:(50,20,2); 7:(-100,-20,0)&$\varnothing$\\
9&128&$(1,23)$&1:(-40,-80,-8); 3:(40,20,-20); 4:(200,80,8); 7:(-100,-20,0)&$\varnothing$\\
10&256&$(2,1)$&1:(-10,-20,-2); 2:(-40,80,0); 3:(160,80,-80); 6:(2760,0,480); 7:(-400,-80,0)&$\varnothing$\\
11&768&$(2,3)$&1:(-10,-20,-2); 2:(-40,80,0); 3:(0,80,-112); 6:(200,0,-32); 7:(-400,-80,0)&$\infty$\\
12&256&$(4,5)$&1:(-10,-20,-2); 2:(-40,80,0); 3:(160,80,-80); 6:(690,0,120); 7:(-400,-80,0)&$\varnothing$\\
\bottomrule\end{longtable}
\endgroup

